\chapter{Tổng quan}

\section{Giới thiệu đề tài}

Thương mại điện tử đang phát triển mạnh mẽ tại Việt Nam với tốc độ tăng trưởng
trung bình 25\% mỗi năm, đưa Việt Nam trở thành một trong những thị trường
thương mại điện tử phát triển nhanh nhất Đông Nam Á. Sự bùng nổ của các nền
tảng như Shopee, Lazada, Tiki đã tạo điều kiện cho hàng triệu người bán nhỏ lẻ
tham gia kinh doanh trực tuyến. Tuy nhiên, một trong những rào cản lớn nhất đối
với người bán là việc soạn mô tả sản phẩm chất lượng --- công việc đòi hỏi
kỹ năng viết quảng cáo, hiểu biết về SEO và tốn nhiều thời gian.

Theo khảo sát, người bán trung bình mất 10--15 phút để soạn mô tả cho một sản
phẩm, bao gồm chụp ảnh, viết tiêu đề, mô tả chi tiết và chọn danh mục phù hợp.
Với những người bán có hàng trăm sản phẩm, đây là một công việc tốn kém đáng kể.
Đặc biệt, nhiều người bán nhỏ lẻ tại các vùng nông thôn không thành thạo việc
soạn thảo văn bản trên máy tính, nhưng lại có khả năng mô tả sản phẩm bằng
lời nói một cách tự nhiên.

Dự án ShopNexus ra đời nhằm giải quyết vấn đề trên bằng cách xây dựng hệ thống
nhúng hỗ trợ bán hàng thông minh, cho phép người bán chỉ cần mô tả sản phẩm
bằng giọng nói thông qua thiết bị ESP32-CAM. Hệ thống tự động thực hiện toàn
bộ quy trình: khử nhiễu âm thanh, nhận dạng giọng nói, sinh mô tả sản phẩm
chuyên nghiệp theo nhiều phong cách viết, và phân loại vào danh mục phù hợp
--- tất cả chỉ với một thao tác nói.

\section{Tình hình nghiên cứu}

Bài toán hỗ trợ bán hàng bằng giọng nói liên quan đến nhiều lĩnh vực
nghiên cứu đang được phát triển mạnh mẽ.

Trong lĩnh vực khử nhiễu giọng nói, các phương pháp dựa trên deep learning
đã cho kết quả vượt trội so với các phương pháp truyền thống như spectral
subtraction hay Wiener filtering. Mô hình DTLN (Dual-signal Transformation
LSTM Network) được đề xuất bởi Westhausen và Meyer (2020) là một trong những
mô hình nhỏ gọn nhất ($\approx$ 1,8 triệu tham số) mà vẫn đạt hiệu năng
real-time, phù hợp cho triển khai trên thiết bị nhúng.

Về nhận dạng giọng nói tiếng Việt, các mô hình như PhoWhisper (VinAI) và
OpenAI Whisper đã đạt được độ chính xác đáng kể. Tuy nhiên, việc triển khai
local trên thiết bị nhúng vẫn là thách thức do yêu cầu tài nguyên tính toán.
Giải pháp local như PhoWhisper-large (VinAI) cung cấp chất lượng cao cho
tiếng Việt mà không phụ thuộc cloud API, tuy nhiên yêu cầu tài nguyên
tính toán đáng kể (khoảng 2 GB VRAM).

Trong phân loại sản phẩm, các mô hình transformer đa ngôn ngữ như XLM-RoBERTa
cho phép phân loại text ở nhiều ngôn ngữ khác nhau mà không cần dữ liệu huấn
luyện riêng cho từng ngôn ngữ, nhờ cơ chế cross-lingual transfer learning.

Điểm khác biệt của ShopNexus so với các nghiên cứu hiện có là việc tích hợp
toàn bộ pipeline --- từ thu thập dữ liệu trên thiết bị nhúng đến xử lý AI
và hiển thị kết quả --- thành một hệ thống end-to-end hoàn chỉnh, tập trung
vào ngữ cảnh thương mại điện tử Việt Nam.

\section{Mục tiêu đề tài}

Đề tài hướng đến các mục tiêu cụ thể sau:

\begin{itemize}
\item Xây dựng thiết bị nhúng ESP32-CAM có khả năng thu âm thanh và hình ảnh
  sản phẩm đồng thời, truyền dữ liệu real-time đến server xử lý thông qua
  giao thức TCP, đảm bảo không mất mát dữ liệu trong quá trình truyền.
\item Phát triển dịch vụ khử nhiễu giọng nói sử dụng mô hình DTLN,
  huấn luyện trên bộ dữ liệu tiếng Việt VIVOS kết hợp với dữ liệu nhiễu
  Microsoft DNS, đạt latency dưới 50 ms để đảm bảo xử lý real-time.
\item Xây dựng dịch vụ nhận dạng giọng nói (Speech-to-Text) và tự động
  sinh mô tả sản phẩm theo nhiều phong cách viết khác nhau, xử lý tốt
  các đặc thù của tiếng Việt đa vùng miền.
\item Phát triển dịch vụ phân loại sản phẩm vào 32 danh mục thương mại
  điện tử sử dụng mô hình XLM-RoBERTa, đạt accuracy trên 85\%.
\item Tích hợp toàn bộ pipeline thành hệ thống end-to-end hoạt động liền mạch
  với giao diện web trực quan, cho phép người dùng thao tác dễ dàng.
\end{itemize}

\section{Phạm vi đề tài}

\begin{itemize}
\item \textbf{Phần cứng:} Thiết bị nhúng ESP32-CAM với microphone INMP441
  (thu âm 16 kHz, mono) và camera OV5640 (JPEG, hỗ trợ đến 1600$\times$1200).
  Kết nối WiFi 802.11 b/g/n, truyền dữ liệu qua TCP.
\item \textbf{Phần mềm:} Server trung gian Node.js (relay TCP$\leftrightarrow$WebSocket),
  dịch vụ khử nhiễu (Flask + TensorFlow), dịch vụ nhận dạng giọng nói
  (FastAPI + OpenAI), dịch vụ phân loại (FastAPI + PyTorch), và giao diện web.
\item \textbf{Ngôn ngữ:} Tập trung xử lý tiếng Việt (đa vùng miền).
  Mô hình phân loại hỗ trợ đa ngôn ngữ nhờ XLM-RoBERTa.
\item \textbf{Danh mục:} 32 danh mục thương mại điện tử phổ biến,
  từ Electronics đến Fashion, Health, và Industrial.
\item \textbf{Giới hạn:} Hệ thống yêu cầu kết nối WiFi ổn định.
  Module STT sử dụng cloud API (OpenAI), chưa hỗ trợ xử lý offline.
\end{itemize}

\section{Cấu trúc báo cáo}

Báo cáo được tổ chức thành 6 chương:

\begin{itemize}
\item \textbf{Chương 1 -- Tổng quan:} Giới thiệu đề tài, tình hình nghiên cứu,
  mục tiêu và phạm vi.
\item \textbf{Chương 2 -- Cơ sở lý thuyết:} Trình bày các kiến thức nền tảng
  về phần cứng ESP32, giao thức truyền thông, và các mô hình AI được sử dụng
  (DTLN, XLM-RoBERTa, GPT-4o).
\item \textbf{Chương 3 -- Thiết kế hệ thống:} Kiến trúc tổng thể,
  luồng xử lý, và thiết kế chi tiết từng thành phần.
\item \textbf{Chương 4 -- Cài đặt hệ thống:} Chi tiết cài đặt phần cứng,
  cấu hình các dịch vụ, và triển khai pipeline xử lý.
\item \textbf{Chương 5 -- Kết quả và đánh giá:} Kết quả đạt được,
  phân tích ưu điểm và hạn chế của hệ thống.
\item \textbf{Chương 6 -- Kết luận và hướng phát triển:} Tổng kết
  và đề xuất hướng phát triển tương lai.
\end{itemize}
